\subsection{信号处理的应用场景与方法体系}

信号处理理论的价值体现在解决实际工程问题中。本节从移动计算领域的主流应用场景出发，系统梳理各类信号处理方法的选择逻辑、实现细节与性能权衡。我们将看到，同一数学工具在不同场景下如何被创造性地组合使用，以及为什么某种方法在特定条件下优于其他选择。

\subsubsection{WiFi感知：从通信到传感}

WiFi信号最初设计用于数据传输，但研究者发现它携带了丰富的环境信息。当人在WiFi链路附近移动时，身体反射和遮挡会改变信号的传播路径，这些变化编码在信道状态信息（CSI）中。CSI不仅包含信号强度（RSSI），还记录了每个子载波的幅度和相位——这构成了WiFi感知的技术基础。

\begin{figure}[htbp]
    \centering
    \begin{mmd}{0.9\textwidth}
flowchart TD
    A[WiFi发射端] -->|原始信号| B[信号预处理]
    B --> C[去噪与清洗]
    C --> D[特征提取]
    D --> E{应用场景}
    E -->|人体活动| F[HAR分类器]
    E -->|手势识别| G[Gesture识别]
    E -->|室内定位| H[定位算法]
    E -->|呼吸监测| I[生命体征检测]
    
    style A fill:#e1f5fe
    style B fill:#fff3e0
    style D fill:#e8f5e9
    style E fill:#fce4ec
    \end{mmd}
    \caption{WiFi感知系统的通用处理流程}
    \label{fig:wifi-sensing-pipeline}
\end{figure}

\textbf{信号预处理的核心挑战}在于CSI数据的特殊性质。商用WiFi设备采集的CSI受限于硬件不完美性：载波频率偏移（CFO）使相位随时间漂移，采样频率偏移（SFO）导致相位与子载波索引成线性关系，还有随机包检测延迟引入的相位跳变。直接使用原始相位几乎不可能获得稳定的环境感知。

标准的相位预处理流程如下。首先进行\textbf{相位解缠绕}（Phase Unwrapping），将限制在 $[-\pi, \pi]$ 的相位展开为连续值。然后执行\textbf{线性去趋势}：对所有子载波相位做线性拟合，从原始相位中减去拟合直线，这同时消除了CFO和SFO的影响。对于多天线系统，\textbf{共轭相乘}（Conjugate Multiplication of CSI）是一种更鲁棒的方法——将两根天线接收的CSI共轭相乘，公共的相位偏移被抵消，只剩下与角度相关的相位差。

\begin{figure}[htbp]
    \centering
    \begin{mmd}{0.85\textwidth}
graph LR
    subgraph 原始CSI问题
        A1[相位跳变] --> B1[解缠绕]
        A2[线性漂移] --> B2[拟合去趋势]
        A3[硬件噪声] --> B3[Hampel滤波]
    end
    
    subgraph 处理方法
        B1 --> C[连续相位]
        B2 --> C
        B3 --> C
    end
    
    C --> D[共轭相乘]
    D --> E[稳定相位差]
    
    style C fill:#e8f5e9
    style E fill:#c8e6c9
    \end{mmd}
    \caption{CSI相位预处理流程与问题-方法映射}
    \label{fig:csi-phase-processing}
\end{figure}

\subsubsection{相位测距原理与技术演进}

相位与距离的物理关系源于电磁波传播的基本方程。信号传播距离 $d$ 引入的相位延迟为
\begin{equation}
    \phi = \frac{2\pi d}{\lambda} = \frac{2\pi f d}{c}
\end{equation}
其中 $\lambda$ 是波长，$f$ 是频率，$c$ 是光速。这个关系看似简单，实则蕴含深刻的工程权衡：相位以 $2\pi$ 为周期，意味着测量存在\textbf{模糊性}——距离 $d$ 和 $d + n\lambda$ 产生相同的相位读数。

\begin{figure}[htbp]
    \centering
    \begin{mmd}{0.9\textwidth}
graph LR
    subgraph 相位测距核心关系
        A[距离 d] -->|φ = 2πd/λ| B[相位 φ]
        B -->|解算| C[距离估计]
    end
    
    subgraph 周期模糊性
        D[真实距离] --> E{相位测量}
        E -->|φ| F[候选距离1: d]
        E -->|φ + 2π| G[候选距离2: d + λ]
        E -->|φ + 4π| H[候选距离3: d + 2λ]
    end
    
    style B fill:#ffeb3b
    style E fill:#ffccbc
    \end{mmd}
    \caption{相位测距的基本关系与周期模糊性问题}
    \label{fig:phase-ranging-principle}
\end{figure}

\textbf{单频测距的局限}显而易见。以2.4GHz WiFi为例，波长 $\lambda \approx 12.5$ cm，这意味着仅凭单频相位无法区分12.5cm整数倍的距离。解决方案之一是\textbf{多频相位融合}。Chronos系统利用WiFi在多个频道上的相位测量，构建超定方程组来解算真实距离。

具体而言，假设在 $K$ 个不同频率上测量相位 $\phi_k$，则
\begin{equation}
    \phi_k = \frac{2\pi f_k d}{c} \mod 2\pi, \quad k = 1, 2, ..., K
\end{equation}

这可以转化为一个整数规划问题：寻找距离 $d$ 和整数 $n_k$，使得
\begin{equation}
    \frac{\phi_k + 2\pi n_k}{2\pi f_k} = \frac{d}{c}, \quad \forall k
\end{equation}

当频率间隔足够大且测量足够精确时，中国剩余定理保证了唯一解的存在。Chronos实现了分米级的测距精度，远超RSSI方法的米级精度。

\begin{figure}[htbp]
    \centering
    \begin{mmd}{0.9\textwidth}
graph TD
    A[多频相位测量] --> B{频率选择}
    B -->|2.4GHz频段| C[多个20MHz信道]
    B -->|5GHz频段| D[更多可用信道]
    
    C --> E[构建方程组]
    D --> E
    
    E --> F[中国剩余定理]
    F --> G[解算整数周期数]
    G --> H[精确距离估计]
    
    subgraph 精度提升
        I[单频精度: ~6cm] --> J[多频融合]
        J --> K[Chronos精度: ~10cm]
    end
    
    style H fill:#c8e6c9
    style K fill:#a5d6a7
    \end{mmd}
    \caption{Chronos多频相位测距原理}
    \label{fig:chronos-ranging}
\end{figure}

\textbf{到达时间（ToF）测距}是另一种基于信号传播时间的方法。ToF与距离的关系为 $d = c \cdot \tau$，其中 $\tau$ 是传播时间。然而，直接测量ToF面临巨大挑战：WiFi信号传播1米仅需3.3纳秒，这要求时钟同步精度达到亚纳秒级，这在商用设备中几乎不可能实现。

Widar和类似系统的创新在于\textbf{消除对绝对时间的依赖}。通过测量两个天线接收同一信号的相位差，可以推导出相对距离差。假设两天线间距为 $d_{ant}$，信号到达角度为 $\theta$，则路径差为 $d_{ant} \sin\theta$，对应的相位差为
\begin{equation}
    \Delta\phi = \frac{2\pi d_{ant} \sin\theta}{\lambda}
\end{equation}

由此可解算到达角
\begin{equation}
    \theta = \arcsin\left(\frac{\lambda \Delta\phi}{2\pi d_{ant}}\right)
\end{equation}

结合多个AP的AoA测量，可以通过三角定位确定目标位置。

\begin{figure}[htbp]
    \centering
    \begin{mmd}{0.85\textwidth}
graph LR
    subgraph 天线阵列几何
        A[发射源] -->|距离r| B[天线1]
        A -->|距离r+Δd| C[天线2]
        B ---|d_ant| C
    end
    
    subgraph 相位差计算
        D[路径差 Δd] -->|Δd = d_ant·sinθ| E[相位差 Δφ]
        E -->|Δφ = 2πΔd/λ| F[到达角θ]
    end
    
    subgraph 定位
        G[AP1的AoA] --> H[三角定位]
        I[AP2的AoA] --> H
        H --> J[目标位置]
    end
    
    style F fill:#e3f2fd
    style J fill:#c8e6c9
    \end{mmd}
    \caption{基于相位差的AoA测距与定位原理}
    \label{fig:aoa-ranging}
\end{figure}

\subsubsection{信号相似度度量：从距离到决策}

信号相似度度量是模式识别的核心问题。在WiFi感知、人体活动识别、室内定位等应用中，我们需要比较两个信号片段是否"相似"，这一决策直接影响系统性能。

\begin{figure}[htbp]
    \centering
    \begin{mmd}{0.95\textwidth}
graph TD
    A[信号相似度度量] --> B[距离度量]
    A --> C[相关性度量]
    A --> D[时间对齐]
    A --> E[分布度量]
    
    B --> B1[欧几里得距离]
    B --> B2[曼哈顿距离]
    B --> B3[切比雪夫距离]
    B --> B4[汉明距离]
    
    C --> C1[皮尔逊相关系数]
    C --> C2[余弦相似度]
    C --> C3[斯皮尔曼相关]
    
    D --> D1[DTW动态时间规整]
    D --> D2[约束DTW]
    D --> D3[FastDTW]
    
    E --> E1[Wasserstein距离]
    E --> E2[MMD最大均值差异]
    E --> E3[KL散度]
    
    style B1 fill:#e3f2fd
    style C1 fill:#e8f5e9
    style D1 fill:#fff3e0
    style E1 fill:#fce4ec
    \end{mmd}
    \caption{信号相似度度量方法分类体系}
    \label{fig:similarity-methods}
\end{figure}

\textbf{欧几里得距离}是最直观的相似度度量。对于两个信号向量 $\mathbf{x}$ 和 $\mathbf{y}$，其定义为
\begin{equation}
    d_{Euclidean} = \sqrt{\sum_{i=1}^{n}(x_i - y_i)^2}
\end{equation}

在室内定位的指纹匹配中，欧几里得距离是首选方法。假设我们在训练阶段采集了各参考点的CSI指纹，在线定位时计算当前测量与所有指纹的欧几里得距离，选择距离最小的 $k$ 个指纹进行加权平均（WKNN算法），即可估计当前位置。

然而，欧几里得距离对信号的绝对幅度敏感。如果两次测量中发射功率不同，或者环境反射条件变化，即使同一位置的信号也可能因幅度差异而被误判为不相似。\textbf{余弦相似度}解决了这一问题
\begin{equation}
    \cos(\theta) = \frac{\mathbf{x} \cdot \mathbf{y}}{||\mathbf{x}|| \cdot ||\mathbf{y}||}
\end{equation}

余弦相似度只关注信号向量的方向，忽略其长度。这在高维CSI数据中特别有用——不同次测量中总能量可能变化，但各子载波间的相对关系保持稳定。

\begin{figure}[htbp]
    \centering
    \begin{mmd}{0.8\textwidth}
graph LR
    subgraph 欧几里得距离
        A1[信号A] --> C1{计算}
        B1[信号B] --> C1
        C1 -->|√Σ(xi-yi)²| D1[距离值]
        D1 -->|受幅度影响| E1[变化敏感]
    end
    
    subgraph 余弦相似度
        A2[信号A] --> C2{计算}
        B2[信号B] --> C2
        C2 -->|A·B/|A||B|| D2[夹角余弦]
        D2 -->|仅关注方向| E2[幅度鲁棒]
    end
    
    style E1 fill:#ffccbc
    style E2 fill:#c8e6c9
    \end{mmd}
    \caption{欧几里得距离与余弦相似度的本质区别}
    \label{fig:euclidean-vs-cosine}
\end{figure}

当信号存在时间偏移时，欧几里得距离和余弦相似度都会失效。考虑步态识别：同一个人以不同速度行走，步态周期不同，即使波形形状相似，简单对齐后计算的距离也会很大。\textbf{动态时间规整（DTW）}正是为解决这个问题而生。

DTW通过动态规划寻找两个时间序列之间的最优对齐路径。设两个序列分别为 $X = (x_1, ..., x_n)$ 和 $Y = (y_1, ..., y_m)$，构建 $n \times m$ 的累积距离矩阵
\begin{equation}
    D(i,j) = d(x_i, y_j) + \min\begin{cases} D(i-1, j) \\ D(i, j-1) \\ D(i-1, j-1) \end{cases}
\end{equation}
其中 $d(x_i, y_j)$ 是局部距离（通常为欧几里得距离）。$D(n,m)$ 即为DTW距离，表示最优对齐路径的累积距离。

\begin{figure}[htbp]
    \centering
    \begin{mmd}{0.9\textwidth}
graph TD
    A[序列X: x1,x2,x3,x4] --> B[构建距离矩阵]
    C[序列Y: y1,y2,y3] --> B
    
    B --> D[局部距离计算]
    D --> E[动态规划累积]
    
    E --> F{对齐路径}
    F -->|路径1| G[直接对齐: 距离大]
    F -->|路径2| H[弯曲对齐: 距离小]
    
    H --> I[DTW距离]
    
    subgraph 约束
        J[Warping Window] --> K[限制搜索范围]
        L[Slope Constraint] --> M[防止过度弯曲]
    end
    
    K --> E
    M --> E
    
    style I fill:#c8e6c9
    style H fill:#e8f5e9
    \end{mmd}
    \caption{DTW动态时间规整的核心思想与约束机制}
    \label{fig:dtw-concept}
\end{figure}

原始DTW的计算复杂度为 $O(nm)$，对于长序列可能过于耗时。工程实践中通常添加\textbf{warping window}约束——只允许在对角线附近 $w$ 个单元格内搜索，将复杂度降至 $O(nw)$。FastDTW则采用多尺度逼近策略：先在粗粒度上计算近似对齐，然后逐层细化，在保持精度的同时大幅提升速度。

\textbf{汉明距离}是另一种重要的相似度度量，尤其适用于二进制编码场景。BiCSI算法将CSI指纹编码为二进制序列，通过异或操作快速计算相似度
\begin{equation}
    d_{Hamming} = \sum_{i=1}^{n}(x_i \oplus y_i)
\end{equation}

汉明距离的计算效率极高（仅为位运算），且存储需求大幅降低（从MB级降至KB级）。在资源受限的边缘设备上，这种轻量级方法使实时定位成为可能。

\subsubsection{时频分析方法与选择}

人体活动包含丰富的频率信息：呼吸频率约0.15-0.5 Hz，行走频率约1-2 Hz，手势动作可达5-10 Hz。不同时频分析方法适用于不同类型的活动。

\begin{figure}[htbp]
    \centering
    \begin{mmd}{0.95\textwidth}
graph TD
    A[时频分析方法] --> B[FFT]
    A --> C[STFT]
    A --> D[小波变换]
    A --> E[EMD/HHT]
    
    B -->|适用| B1[平稳信号]
    B -->|优点| B2[计算高效]
    B -->|缺点| B3[无时间定位]
    
    C -->|适用| C1[非平稳信号]
    C -->|优点| C2[时频联合]
    C -->|缺点| C3[分辨率权衡]
    
    D -->|适用| D1[多尺度信号]
    D -->|优点| D2[自适应分辨率]
    D -->|缺点| D3[基函数选择]
    
    E -->|适用| E1[非线性信号]
    E -->|优点| E2[数据驱动]
    E -->|缺点| E3[计算复杂]
    
    style B fill:#e3f2fd
    style C fill:#e8f5e9
    style D fill:#fff3e0
    style E fill:#fce4ec
    \end{mmd}
    \caption{时频分析方法的适用场景与权衡}
    \label{fig:time-frequency-methods}
\end{figure}

\textbf{快速傅里叶变换（FFT）}是频域分析的基础工具。对于离散序列 $x[n]$，其DFT为
\begin{equation}
    X[k] = \sum_{n=0}^{N-1} x[n] e^{-j2\pi kn/N}, \quad k = 0, 1, ..., N-1
\end{equation}

FFT通过蝶形运算将复杂度从 $O(N^2)$ 降至 $O(N\log N)$。在呼吸检测中，对30秒CSI数据做FFT，在0.15-0.5 Hz范围内寻找峰值，即可准确估计呼吸率。

然而，FFT假设信号在整个分析窗内是平稳的，无法捕捉频率随时间的变化。\textbf{短时傅里叶变换（STFT）}通过滑动窗口解决了这一问题
\begin{equation}
    \text{STFT}(t,f) = \int_{-\infty}^{\infty} x(\tau)w(\tau-t)e^{-j2\pi f\tau}d\tau
\end{equation}
其中 $w(\tau-t)$ 是滑动窗函数（通常为Hanning或Hamming窗）。

STFT的核心局限是\textbf{不确定性原理}：时间分辨率和频率分辨率无法同时达到最优。窗口越长，频率分辨率越高，但时间定位越模糊；窗口越短，时间定位越精确，但频率分辨率下降。

\textbf{小波变换}提供了更灵活的解决方案。与STFT使用固定频率的基函数不同，小波变换使用缩放和平移的小波基 $\psi_{a,b}(t) = \frac{1}{\sqrt{a}}\psi\left(\frac{t-b}{a}\right)$，其中 $a$ 是尺度参数（对应频率），$b$ 是平移参数（对应时间）。

离散小波变换（DWT）通过多分辨率分析将信号分解为近似系数（低频）和细节系数（高频）。Daubechies小波（db）因其紧凑支撑和正交性而广泛应用——db1（Haar）适合检测边缘，db5适合平滑信号分析。

\begin{figure}[htbp]
    \centering
    \begin{mmd}{0.9\textwidth}
graph TD
    A[原始信号] --> B[DWT分解]
    
    B --> C1[第一层细节系数]
    B --> C2[第一层近似]
    
    C2 --> D1[第二层细节系数]
    C2 --> D2[第二层近似]
    
    D2 --> E1[第三层细节系数]
    D2 --> E2[第三层近似]
    
    subgraph 频率范围
        C1 -->|高频| F1[快速运动]
        D1 -->|中高频| F2[手势]
        E1 -->|中频| F3[行走]
        E2 -->|低频| F4[呼吸]
    end
    
    style A fill:#e3f2fd
    style E2 fill:#c8e6c9
    \end{mmd}
    \caption{DWT多层分解与人体活动频率对应关系}
    \label{fig:dwt-decomposition}
\end{figure}

CARM系统采用12层DWT分解，覆盖0.15 Hz到300 Hz的频率范围，每层能量反映了特定速度范围的人体运动。这种多尺度特征对复杂活动识别至关重要。

\subsubsection{应用场景的方法组合}

实际系统很少使用单一方法，而是根据场景特点组合多种技术。

\begin{figure}[htbp]
    \centering
    \begin{mmd}{0.95\textwidth}
graph TD
    subgraph 人体活动识别HAR
        A1[原始CSI] --> B1[Hampel滤波]
        B1 --> C1[PCA降维]
        C1 --> D1[DWT分解]
        D1 --> E1[小波能量特征]
        E1 --> F1[SVM/CNN分类]
    end
    
    subgraph 室内定位
        A2[CSI幅度] --> B2[去噪]
        B2 --> C2[欧几里得距离]
        C2 --> D2[k-NN匹配]
        D2 --> E2[WKNN定位]
    end
    
    subgraph 手势识别
        A3[CSI流] --> B3[相位处理]
        B3 --> C3[STFT]
        C3 --> D3[频谱图]
        D3 --> E3[CNN特征]
        E3 --> F3[DTW对齐]
    end
    
    subgraph 呼吸检测
        A4[相位] --> B4[Butterworth滤波]
        B4 --> C4[FFT]
        C4 --> D4[峰值检测]
    end
    
    style F1 fill:#c8e6c9
    style E2 fill:#c8e6c9
    style F3 fill:#c8e6c9
    style D4 fill:#c8e6c9
    \end{mmd}
    \caption{不同应用场景的信号处理方法组合}
    \label{fig:application-pipelines}
\end{figure}

\textbf{人体活动识别（HAR）}的典型流程是：Hampel滤波去除异常值 → PCA提取主成分 → DWT分解 → 提取各层小波能量 → 用SVM或CNN分类。这种组合充分利用了各方法的优势：Hampel滤波鲁棒去除脉冲噪声，PCA压缩数据并去除背景，DWT提供多尺度特征，分类器完成最终决策。

\textbf{室内定位}追求实时性和低计算开销。CSI幅度经过简单滤波后，直接与指纹库计算欧几里得距离，选择最近邻或k近邻加权平均。BiCSI算法进一步优化，将指纹编码为二进制序列，用汉明距离实现毫秒级匹配。

\textbf{手势识别}需要捕捉快速变化的动态特征。相位预处理消除硬件漂移后，STFT将时序转换为时频图（spectrogram），CNN自动提取空间特征，DTW处理不同速度的手势变体。

\textbf{呼吸检测}关注极低频信号（<0.5 Hz）。Butterworth低通滤波保留呼吸成分，FFT提取主导频率，简单但有效。更精细的系统可能使用带通滤波器隔离呼吸和心跳信号。

\subsubsection{性能权衡与工程决策}

选择信号处理方法需要在精度、延迟、计算复杂度和鲁棒性之间权衡。

\begin{table}[htbp]
    \centering
    \caption{信号处理方法性能对比}
    \label{tab:method-comparison}
    \begin{tabular}{|l|c|c|c|c|}
        \hline
        \textbf{方法} & \textbf{精度} & \textbf{计算量} & \textbf{延迟} & \textbf{鲁棒性} \\
        \hline
        欧几里得距离 & 中 & 低 & 低 & 低 \\
        余弦相似度 & 中 & 低 & 低 & 中 \\
        DTW & 高 & 高 & 高 & 高 \\
        汉明距离 & 中 & 极低 & 极低 & 中 \\
        FFT & 中 & 低 & 中 & 高 \\
        STFT & 高 & 中 & 中 & 高 \\
        DWT & 高 & 中 & 低 & 高 \\
        CNN特征 & 很高 & 很高 & 中 & 很高 \\
        \hline
    \end{tabular}
\end{table}

资源受限的边缘设备（如ESP32）倾向于选择汉明距离、FFT等轻量级方法；云端服务器可以承担DTW、CNN等计算密集型任务。实际部署中，预处理通常在设备端完成（去噪、降维），特征提取和分类可选择在边缘或云端执行，取决于具体应用的延迟要求。

跨域泛化是WiFi感知面临的长期挑战。同一模型在不同环境、不同设备上的表现可能显著下降。解决方法包括：数据增强（添加噪声、时间拉伸）、域适应（DANN、梯度反转）、以及测试时适应（TTA、连续学习）。这些方法虽然增加了系统复杂度，但对实际部署至关重要。
